\centerline{\bf \Huge Математический бой. Старшая группа}

\begin{flushright}
        \scriptsize{Только неудачники просят о смерти, проиграв бой... \\Вместо того, чтобы принять поражение, после которого ты должен был умереть!}
\end{flushright}

\problem{1}
 На летних сборах в ЭлМШ Герман решил не решать геометрию вовсе. Баир же решил решать только геометрию. Вскоре РВ посмотрел на рейтинговую таблицу и заметил, что модуль разности их баллов кратен $11$. Какое количество задач они не могли решить в сумме даже в теории, если в каждом листочке от $0$ до $9$. (Приведите все такие варианты и объясните, почему других нет.)

\noindent
\textbf{Ответ.} $1, 3, 5, 7, 9.$ \textbf{Решение.} Любую чётную сумму цифр дают числа вида $1\ldots1$, сумму $11$ --- число $209$, любую нечётную, начиная с $13$, числа вида $1\ldots 1209.$ Пусть число $B$ имеет нечётную сумму цифр, не большую $9$. Тогда знакопеременная сумма его цифр также по модулю не больше 9, и должна делиться на $11$, если число делится на $11$. Но в таком случае она равна $0$, что невозможно, так как она нечётна.


\problem{2}
 Действительные числа $a, b, c, d$ таковы, что $a+b = cd$ и $c+d = ab.$ Докажите неравенство $(a+1)(b+1)(c+1)(d+1) \geq 0$.

\noindent
\textbf{Решение.}
($a+1)(b+1) = ab+a+b+1 = a+b+c+d+1.$ Аналогично, $(c+1)(d+1) = a+b+c+d+1.$ Поэтому $(a+1)(b+1)(c+1)(d+1) = (a+b+c+d+1)^2 \geq 0.$


\problem{3}
Дан квадрат со стороной, большей $1000$. Докажите, что из него можно вырезать не более $4$ клеток таким образом, что остаток можно разрезать на прямоугольники $1\times4$.

\noindent
\textbf{Решение.}
У квадрата со стороной $4n+1$ достаточно вырезать угловую клетку: оставшуюся часть можно разрезать на квадрат со стороной $4n$ и две полоски $1\times4n$, которые без труда режутся на прямоугольники $1\times4.$ В квадрате со стороной $4n+3$ выделим угловой квадрат $7\times7$ и вырежем в нем центральную клетку. Оставшуюся часть квадрата $7\times7$ нетрудно разрезать на четыре прямоугольника $4\times3$, а оставшаяся часть квадрата со стороной $4n+3$ режется на квадрат со стороной $4n-$4 и два прямоугольника $7\times(4n-4)$, легко режущиеся на прямоугольники $1\times4.$

\problem{4}
Дан треугольник ABC, где $\angle A = 70^{\circ}, \angle B = 30^{\circ}$ и $\angle C = 80^{\circ}$. Точка $D$ внутри треугольника такова, что $\angle DAB = 30^{\circ}$ и $\angle DBA = 10^{\circ}.$ Найдите угол $\angle DCB.$

\noindent \textbf{Ответ.} $50^{\circ}$. \textbf{Решение.} Пусть $O$ --- такая точка на отрезке $BD$, что $\angle OAB = 10^{\circ}.$ Тогда треугольник $AOB$ --- равнобедренный.

\newpage

\hspace{3mm}

С другой стороны, $\angle AOB = 160^{\circ} = 2\angle ACB.$ Следовательно, $O$ --- центр описанной окружности треугольника $ABC$. Таким образом, $OA = OC$, а также $\angle CAO = \angle CAB-\angle OAB = 60^{\circ}$. Значит, треугольник $OAC$ --- равносторонний. Наконец, $\angle DAO = \angle DAB-\angle OAB = 30^{\circ}-10^{\circ} = 20^{\circ}, \angle DOA$ --- внешний к треугольнику $AOB$, т.е. $\angle DOA = 2\angle OBA = 20^{\circ} = \angle DAO,$ откуда $DA = DO.$ Следовательно, $D$ лежит на серединном перпендикуляре к отрезку $AO$, и потому $CD$ --- биссектриса угла $ACO$. Значит, $\angle DCB = \angle ACB-\angle ACD = 80^{\circ}-30^{\circ} = 50^{\circ}.$


\problem{5}
На сборах ЭлМШ $100$ учеников, каждый из которых может быть Киви, Лососем или Шиншиллой (Шиншиллы говорят правду, Лососи всегда врут, а Киви могут говорить как правду, так и ложь). $50$ участников сборов сказали: «На сборах Шиншилл больше, чем Киви», а $50$ остальных: «На сборах Лососей больше, чем Киви». Докажите, что Киви на сборах не меньше $25.$

\textbf{Решение.} Докажем от противного. Пусть Киви меньше 25. Заметим, что Лососей и Шиншилл не может быть одновременно больше 25, иначе оба утверждения будут правдивыми и Лососям просто нечего отвечать. Предположим, что Лососей меньше 25. Тогда если учеников сумммарно 100, а Лососей и Киви меньше 25, то Шиншилл хотя бы 51. Но учитывая, что Шиншиллы всегда говорят одно и то же, то и одинаковых ответов должно быть хотя бы 51, что не верно. Противоречие.


\problem{6}
Будем называть точку плоскости узлом, если обе её координаты --- целые числа. Внутри некоторого треугольника с вершинами в узлах лежит ровно два узла (возможно, какие-то еще узлы лежат на его сторонах). Докажите, что прямая, проходящая через эти два узла, либо проходит через одну из вершин треугольника, либо параллельна одной из его сторон.

\textbf{Решение.}  Если прямая $XY$ не проходит через вершины треугольника $ABC$, то она не пересекает одну из его сторон, например, $AB$. Тогда внутри треугольников $AXB$ и $AYB$ нет узлов, и по формуле Пика площади треугольников $AXB$ и $AYB$ равны. Значит, точки $X$ и $Y$ находятся на равных расстояниях от прямой $AB$, то есть прямая $XY$ параллельна прямой $AB$.

%https://problems.ru/view_problem_details_new.php?id=32889


\problem{7}
6 идущих подряд натуральных чисел каким-то образом разбивают на две тройки чисел. Затем считают произведения чисел в этих тройках и из большего произведения вычитают меньшее. Какое наименьшее значение может принимать результат?

\newpage

\hspace{3mm}

\textbf{Ответ.} $2$. \textbf{Решение.} Оценка. Допустим, разность двух произведений равна $1$. Тогда они имеют разную чётность, и все чётные числа из шестёрки должны составлять одну тройку, а все нечётные $\dfrac{3}{4}$ другую. Но нетрудно проверить, что разность $(n+1)(n+3)(n+5)-n(n+2)(n+4)$ больше $1$. Допустим, два произведения равны. Заметим, что среди трёх идущих подряд нечётных чисел есть число $k$, не делящееся ни на $3$, ни на $5$. Если это единица, то у нас числа от $1$ до $6$, и произведение, содержащее пятёрку, не может равняться другому. Если же это не единица, то все простые делители числа $k$ больше $5$. Значит, $k$ взаимно просто с остальными числами шестёрки, и потому произведение, содержащее $k$, не может равняться второму произведению. Пример. $3\times 4 \times 6 - 2 \times 5 \times 7 = 2.$

\problem{8}
В треугольнике ABC проведена биссектриса $AL$. Точка $P$ --- основание перпендикуляра, опущенного из точки $B$ на отрезок $AL$. Из точки $P$ на сторону $AB$ опущена высота $PQ$. На стороне AB выбрана такая точка $R$, что $RQ = AC/4.$ Докажите, что $AB+BC \geq 4PR.$

\noindent\textbf{Решение.} Опустим из $P$ перпендикуляр $PS$ на прямую $AC$. Отложим отрезок $ST = RQ$ на $AC$ (не важно, в какую сторону). По свойству биссектрисы мы имеем $PS = PQ.$ Треугольники $PQR$ и $PST$ равны по признаку равенства прямоугольных треугольников, откуда $PR = PT$. Остается доказать, что $AB+BC \geq 4PT.$

Продлим перпендикуляр $BP$ до пересечения с прямой $AC$ в точке $E$. В треугольнике $ABE$ отрезок $AP$ --- высота и биссектриса, а значит и медиана. Проведем среднюю линию $A_1C_1$ треугольника $ABC$. Она делит пополам любой отрезок, соединяющий $B$ c точкой на стороне $AC$. Поэтому P лежит на прямой $A_1C_1$.

Пусть $M$ --- середина $AC$. Восставим перпендикуляр $MB'$ к прямой $AC$, длина которого равна высоте $BH$ треугольника $ABC$, в сторону точки $B$. Средняя линия $A1C1$ делит $B'M$ пополам и параллельна AC, поэтому $B'M = 2PS.$ Кроме того, $AM = MC = 2ST$, поэтому $B'C = B'A = 2PT.$

Отметим на продолжении отрезка $AB'$ за $B'$ точку $C_2$ так, что $C_2B' = AB'.$ Пусть $l$ --- это прямая, проходящая через $B'$ параллельно $AC$. Тогда $B$ лежит на $l$, а точки $C_2$ и $C$ симметричны относительно прямой $l$. Поэтому $C_2B = BC.$ Следовательно, $AB+BC = AB+BC_2 \geq AC2 = AB'+B'C_2 = 4PT = 4RS$, что и требовалось доказать.